% Canonical Introduction; final score definition and interpretation pass, 5 October 2026.
\section*{Introduction}

Drug-induced liver injury is the most common cause of acute liver failure in the Western world,\cite{ref001,ref002} and acetaminophen (APAP) overdose in particular is the leading cause of acute liver failure in the United States and the United Kingdom.\cite{ref003,ref004} Mechanistically, APAP hepatotoxicity produces a characteristic pattern of centrilobular, pericentral hepatocyte necrosis, reflecting the high CYP2E1-dependent bioactivation of APAP in zone-3 hepatocytes.\cite{ref004,ref005} Surviving hepatocytes can subsequently follow heterogeneous trajectories, including restoration toward baseline hepatocyte states as the liver regenerates,\cite{ref006,ref007} whereas a subset can remain in a state of persistent compensation for days after injury. What determines whether an individual hepatocyte resolves toward recovery or remains in a compensatory or injury-associated state is still incompletely understood at single-cell resolution. Resolving this heterogeneity is relevant to understanding clinical recovery and regenerative failure in acute liver injury, since the adequacy of hepatocyte regeneration is a key determinant of outcome after APAP-induced acute liver failure.\cite{ref006,ref008}

Single-cell RNA sequencing (scRNA-seq) enables transcriptional profiling of thousands of hepatocytes across injury and recovery,\cite{ref009,ref010,ref011} and trajectory inference (TI) methods order cells along a pseudotime axis relative to a reference state.\cite{ref012,ref013} Expression similarity alone can be difficult to interpret when injury-associated cells reacquire metabolic programs during recovery or persistent compensation.\cite{ref010} External directionality from RNA velocity\cite{ref014,ref015} or lineage tracing\cite{ref016,ref017} requires paired count layers or clonal barcodes absent from many deposited datasets, including two of the four integrated here. Controlled injury studies also provide sampling time and treatment annotations that can constrain transition graphs. These annotations distinguish exposure history even when treated and untreated hepatocytes have similar transcriptional profiles; they do not, by themselves, establish an individual cell's fate. Injury and recovery involve temporally distinct responses,\cite{ref006,ref018,ref019} motivating treatment-aware modeling of the drug- and toxicant-response time courses now available through scRNA-seq.\cite{ref020,ref021,ref022}

Graph neural networks (GNNs) propagate information over single-cell graphs to represent cell states\cite{ref023,ref024} and can complement per-cell transcriptional scoring of disease-associated programs.\cite{ref025,ref026,ref027} A graph shared across an entire dataset may obscure changes in transition structure between response phases. Multi-scale and multi-resolution architectures provide ways to represent scale-dependent structure outside biology\cite{ref028,ref029,ref030,ref031} and in medical time series.\cite{ref032} APAP responses progress through early stress, peak injury, heterogeneous recovery or compensation, and later resolution,\cite{ref006,ref033,ref034} within the broader sequence of inflammation, proliferation, and remodeling in tissue repair.\cite{ref035,ref036,ref037} Phase-specific graphs may therefore reveal structure that a single propagation scale does not resolve. Network-science analyses of directed edges\cite{ref038,ref039} and cell-state transitions\cite{ref040,ref041} offer a complementary way to examine where model-derived information concentrates within these graphs.

Here we present TrajRisk, a framework that learns directed, treatment-aware transition probabilities at multiple temporal scales and integrates expression-derived node features across these graphs via a graph neural network. The TrajRisk score is a dimensionless model estimate of normalized Monocle 3 pseudotime in the mouse injury--recovery reference; ``risk'' is the method's operational label for this trajectory-state score and does not estimate a probability of future injury, failure to recover, or a clinical outcome. The score depends on a cell's expression and its incoming transition neighborhood. TrajRisk trains a Random Forest\cite{ref042} for each temporal scale using nearest-neighbor-defined cell-pair labels and features that include sampling time, treatment, and pseudotime. Parallel convolution branches apply a directed operation suited to asymmetric transitions\cite{ref043} rather than symmetric GCN normalization,\cite{ref044} then fuse the scale-specific embeddings. We evaluate the framework across four integrated mouse APAP studies (36,067 hepatocytes; 6,582 reference cells across five sampling times). TrajRisk improves average precision over expression-only XGBoost within the 24~h validation subset, whereas XGBoost performs better on global AUROC and at later recovery timepoints. External mouse comparisons show cohort-dependent performance. After mouse-only retraining in shared mouse--human coordinates, TrajRisk scores distinguish APAP-ALF from healthy donors without using human condition labels for supervised model training.\cite{ref019} We also examine orthologous gene signatures in human nuclei and spatial tissue, extending cross-species transcriptional comparisons\cite{ref045,ref046} while distinguishing positive findings from null results. Finally, flux analysis identifies a candidate 24--48~h resolution gateway for experimental testing. These analyses evaluate trajectory-state representation and transfer, rather than clinical prognosis or causal recovery mechanisms.
